\documentclass[../main.tex]{subfiles}
\begin{document}

\section{Methods}
\label{sec:methods}

\subsection{Dataset}
\label{sec:dataset}

The evaluation corpus consists of randomised controlled trials (RCTs) included
in Cochrane intervention systematic reviews published in 2024.
Beginning from a source list of 149 reviews, RCTs were selected through a
stratified, risk-of-bias-aware sampling procedure capped at three RCTs per
review, and a trial was retained only when a full-text report (PDF) of the
trial could be obtained.
This procedure yielded a corpus of \textbf{227 RCTs} drawn from 123 distinct
Cochrane reviews (trial publication years 1966--2023).
The full construction pipeline and the resulting labelled dataset are described
in a companion data paper \citep{Dataset2026}.

Risk of bias annotations were obtained directly from the published Cochrane
review reports and reflect the consensus judgements produced by the original
review teams following standard Cochrane procedures (independent assessment by
at least two trained reviewers, with disagreements resolved by discussion or
arbitration).
These annotations were not re-produced or independently adjudicated for the
present study; they represent the published Cochrane consensus, which constitutes
the reference standard used by the field.
Although the revised Risk of Bias~2.0 tool (RoB~2; \citealt{Sterne2019}) is
now recommended for new systematic reviews, it reorganises the domain structure
and changes the rating scale, rendering it incompatible with the existing
RoB~1.0 corpus.
RoB~1.0 remains the basis for the large majority of existing systematic review
annotations --- including all 123 reviews sampled here and most of the large
labelled datasets used for training and evaluation of automated systems --- so
this study's focus on RoB~1.0 is both corpus-driven and practically motivated.
Each study was annotated using the Cochrane Risk of Bias~1.0 (RoB~1.0) tool
\citep{Higgins2011}, which covers seven methodological domains:
D1~(random sequence generation),
D2~(allocation concealment),
D3~(blinding of participants and personnel),
D4~(blinding of outcome assessment),
D5~(incomplete outcome data),
D6~(selective outcome reporting), and
D7~(other sources of bias).
Each domain received one of three judgements---\emph{Low Risk}, \emph{High Risk},
or \emph{Unclear Risk}---so every study contributes seven domain-level cells,
giving 1{,}589 evaluated cells ($227 \times 7$) in total.
The ground-truth class distribution is markedly imbalanced across domains
(reported in Section~\ref{sec:corpus} and Table~\ref{tab:class-dist}), a property
that motivates the agreement metrics adopted in Section~\ref{sec:metrics}.

\textbf{Domain normalisation.}
Not all source reviews reported risk of bias under the seven canonical RoB~1.0
labels: some used a single combined blinding entry (performance and detection
bias) or split individual domains into several sub-domains.
Prior to evaluation, all annotations were normalised to the seven canonical
domains---combined or split blinding entries were mapped to D3~(blinding of
participants and personnel) and D4~(blinding of outcome assessment), and other
custom sub-domains were consolidated into the corresponding canonical domain.
Where merged sub-domains carried conflicting judgements, the more conservative
(higher-risk) judgement was retained.
Normalisation affected 73 of the 227 studies (32\%), predominantly the blinding
domains (D3, D4); the complete mapping is documented in the companion data
paper \citep{Dataset2026}.

\textbf{Intervention type.}
Each study was assigned a single intervention category:
drug/medication, non-pharmacological intervention, procedure, or medical device.

\textbf{Clinical domain.}
Studies were tagged with one or more Cochrane Review Group categories
(multi-label).
For stratified analyses (Section~\ref{sec:methods-strata}), each multi-domain
study was included in every applicable group.

\subsection{Document Preparation Pipeline}
\label{sec:pipeline}

Full-text study reports (PDF) were processed through a two-stage pipeline
that converts each article into a standardised plain-text representation
before it is submitted to any model.
Converting to a common format ensures that all eight models receive identical
textual input, so that observed differences in performance reflect model
behaviour rather than differences in how each model internally parses or
prioritises content from a raw PDF.

\textbf{Stage~1 --- PDF-to-Markdown conversion.}
Each PDF was converted to structured Markdown using MinerU \citep{Wang2024MinerU},
which preserves headings, tables, and inline text whilst replacing embedded figures
with placeholder tokens of the form \texttt{[Figure~$k$]}.

\textbf{Stage~2 --- Figure description.}
Each figure placeholder was resolved by Gemini~2.5~Flash
(\texttt{google/gemini-2.5-flash}, accessed via OpenRouter in May--June~2026),
a vision-language model (VLM) that received the corresponding page image
together with surrounding article context and generated a structured textual
description following a fixed eight-field schema (figure type, axes and units,
legend and groups, visual summary, plotted patterns, and unreadable or uncertain
areas; full prompt reproduced in Supporting Information~S10).
The description was then inserted in place of the placeholder.
After Stage~2, the document was a self-contained Markdown file in which all
information---including any figures present in the source document---was represented as plain text.
The result is a self-contained Markdown file in which the full article
content---text, tables, and figures---is represented uniformly as plain text,
forming the common input submitted to all eight models.

\subsection{Risk of Bias Assessment}
\label{sec:rob-assessment}

All models were evaluated in a zero-shot, single-agent setting: each processed
Markdown document was submitted to the model with a fixed assessment prompt
(reproduced in full in Supporting Information~S1), which specified the domain
criteria, the three admissible judgements (Low Risk, High Risk, Unclear Risk),
and the required JSON output format.
No few-shot examples were provided, and the prompt was used without modification
across all models.

\textbf{Inference modes.}
We evaluated two inference strategies that differ in how the seven RoB domains
are requested from the model.
In the \emph{per-domain} mode, seven sequential inference calls are made per
study: each call supplies the full Markdown text together with a single domain
identifier (e.g.\ \texttt{Domain: D1}), and the model returns a JSON object
containing a judgement and justification for that domain only.
In the \emph{single-call} mode, a single inference call per study is made:
the model receives the full Markdown text and is instructed to return a JSON
array containing judgements and justifications for all seven domains
simultaneously.
All eight models were evaluated under the per-domain mode.
Per-domain mode was adopted as the primary protocol because each RoB~1.0
domain represents an independent methodological criterion; separating calls
ensures that each domain is assessed under a dedicated instruction context,
without the model being required to simultaneously manage seven distinct
evaluation tasks in a single generation pass.
Gemini~2.5~Flash and DeepSeek~v4-Flash were additionally evaluated under the
single-call mode on the same 227 studies, enabling a controlled comparison of
the two inference strategies whilst holding model, prompt content, and corpus
constant.
These two models were selected because their published list prices per token on
OpenRouter were the lowest among the eight models evaluated, making a second
full-corpus run economically practical.
The results of this comparison are reported in Section~\ref{sec:inference}.

\subsection{Language Models}
\label{sec:models}

Eight frontier LLMs were evaluated (Table~\ref{tab:models}), accessed via the
OpenRouter~API in June~2026 using the identifier strings listed therein.
No model-specific fine-tuning, prompt engineering beyond the fixed assessment
prompt (Section~\ref{sec:rob-assessment}), or retrieval augmentation was applied.

\textbf{Evaluation protocol.}
All calls used provider-default generation parameters and each study--condition
pair received a single inference draw with no output aggregation, constituting
a hyperparameter-free single-draw evaluation protocol.
This design reflects two constraints.
First, decoding parameters such as temperature interact with model architecture
and training objective in model-family-specific ways; varying them across models
would confound generation configuration with prompt design, rendering
cross-model comparisons uninterpretable.
Second, post-hoc reliability strategies---including self-consistency decoding
\citep{Wang2023}, output ensembling, and majority voting over repeated
draws---are optimisations orthogonal to native zero-shot capability and
constitute a separate experimental question outside the scope of this benchmark.

\textbf{Reproducibility and versioning.}
OpenRouter routes each identifier string to the provider-current checkpoint at
inference time; underlying weights may be updated by the provider without a
corresponding change to the public identifier.
All results are therefore conditioned on the checkpoint states active during the
June~2026 inference window and should not be assumed to generalise to subsequent
revisions of the same identifier strings.

\begin{table}[ht]
\centering
\caption{Language models evaluated. All models accessed via the OpenRouter API
in June~2026 using the identifier strings shown; zero-shot single-agent
mode throughout; all 227 studies evaluated per model.
Context length: maximum tokens accepted per call as reported by OpenRouter at
access time.
Release date: date the model became available on OpenRouter.
Exact checkpoint versions are opaque at the provider level; the identifiers
shown are the strings passed to the API and represent the model state at the
time of inference.}
\label{tab:models}
\begin{tabular}{llcc}
\toprule
Model & OpenRouter identifier & Context & Release \\
\midrule
Gemini~2.5~Pro    & \texttt{google/gemini-2.5-pro}          & 1M   & Jun 2025 \\
Gemini~2.5~Flash  & \texttt{google/gemini-2.5-flash}        & 1M   & Jun 2025 \\
GPT-5.4           & \texttt{openai/gpt-5.4}                 & 1M   & Mar 2026 \\
Claude~Haiku~4.5  & \texttt{anthropic/claude-haiku-4.5}     & 200K & Oct 2025 \\
DeepSeek~v4-Flash & \texttt{deepseek/deepseek-v4-flash}     & 1M   & Apr 2026 \\
Grok~4.3          & \texttt{x-ai/grok-4.3}                  & 1M   & Apr 2026 \\
Qwen~3.7-Plus     & \texttt{qwen/qwen3.7-plus}              & 1M   & May 2026 \\
GLM-5             & \texttt{z-ai/glm-5}                     & 203K & Feb 2026 \\
\bottomrule
\end{tabular}
\end{table}

\subsection{Evaluation Metrics}
\label{sec:metrics}

\textbf{Accuracy.}
The proportion of study--domain cells in which the model's judgement matched
the ground-truth label, equivalent to the observed agreement $p_o$ used in
the chance-corrected coefficients below.
Let $\mathbf{N} = (N_{ij})$ denote the $K \times K$ confusion matrix (here, $K=3$),
with rows indexed by predicted class $i$ and columns indexed by ground-truth
class $j$:
\begin{equation}
  \text{Accuracy} = p_o = \frac{\sum_{k=1}^{K} N_{kk}}{N}
  = \frac{\sum_{k=1}^{K} N_{kk}}{\sum_{i=1}^{K}\sum_{j=1}^{K} N_{ij}},
  \label{eq:accuracy}
\end{equation}
where $N_{ij}$ is the number of cells with ground-truth class $j$ predicted as
class $i$, and
\(
N = \sum_{i=1}^{K}\sum_{j=1}^{K} N_{ij}
\)
is the total number of evaluated cells ($N = 1{,}589$ overall; $N = 227$ per
domain; smaller for intervention-type and clinical-domain strata).

\textbf{Per-class F1.}
RoB assessment is a three-class problem ($K = 3$).
Per-class metrics are computed under a one-vs-rest scheme: each study--domain cell is
treated as a binary classification (``class $k$'' vs.\ ``not class $k$'') for each class
$k \in \{\text{Low Risk, High Risk, Unclear Risk}\}$,
yielding four counts:
\begin{itemize}
  \item $\mathrm{TP}_k = N_{kk}$: model and ground truth both assigned class $k$;
  \item $\mathrm{FN}_k = \sum\limits_{i \neq k} N_{ik}$: ground truth was $k$, model assigned another class
        (positive missed);
  \item $\mathrm{FP}_k = \sum\limits_{j \neq k} N_{kj}$: model assigned $k$, ground truth was another class
        (false alarm);
  \item $\mathrm{TN}_k = \sum\limits_{i \neq k}\sum\limits_{j \neq k} N_{ij}$: neither model nor ground truth assigned class $k$.
\end{itemize}
Precision (positive predictive value) and recall (sensitivity) for class $k$ are
\begin{equation}
  \mathrm{Precision}_k = \frac{\mathrm{TP}_k}{\mathrm{TP}_k + \mathrm{FP}_k},
  \qquad
  \mathrm{Recall}_k    = \frac{\mathrm{TP}_k}{\mathrm{TP}_k + \mathrm{FN}_k},
  \label{eq:prec-rec}
\end{equation}
and their harmonic mean is the F1 score for class $k$:
\begin{equation}
  \mathrm{F1}_k = \frac{2 \cdot \mathrm{Precision}_k \cdot \mathrm{Recall}_k}
                       {\mathrm{Precision}_k + \mathrm{Recall}_k}.
  \label{eq:f1}
\end{equation}
Per-class F1 reveals class-specific failure modes independently of overall performance.

\textbf{Macro-F1.}
Per-class F1 is computed over the full pool of 1{,}589 cells; the three class-level
F1 scores are then averaged with equal weight:
\begin{equation}
  \text{Macro-F1} = \frac{1}{K}\sum_{k=1}^{K} \mathrm{F1}_k, \qquad K = 3.
  \label{eq:macro-f1}
\end{equation}
Unlike Accuracy, Macro-F1 assigns equal weight to each category regardless of
prevalence, penalising models that achieve high accuracy by concentrating
predictions on the majority class.

\textbf{Cohen's $\kappa$.}
A chance-corrected agreement coefficient that adjusts observed agreement $p_o$
for the agreement expected by chance under the assumption that annotators label
independently according to their observed marginal distributions:
\begin{equation}
  \kappa = \frac{p_o - p_e^\kappa}{1 - p_e^\kappa},
  \qquad
  p_e^\kappa = \sum_{k=1}^{K} \hat{p}_{k,\text{model}} \cdot \hat{p}_{k,\text{gt}},
  \label{eq:kappa}
\end{equation}
where
\(
  \hat{p}_{k,\text{model}} = \dfrac{1}{N}\sum_{j=1}^{K} N_{kj}
\)
and
\(
  \hat{p}_{k,\text{gt}} = \dfrac{1}{N}\sum_{i=1}^{K} N_{ik}
\)
are the marginal proportions of class $k$ in the model predictions and ground
truth, respectively (row and column sums of $\mathbf{N}$ divided by $N$).
$\kappa$ is reported alongside AC1 to expose their divergence under class
imbalance, illustrating the prevalence paradox (see Gwet's AC1 below).

\textbf{Gwet's AC1.}
A chance-corrected agreement coefficient robust to the prevalence paradox
\citep{Gwet2008, Feinstein1990}: when marginal distributions are highly skewed,
classical chance-corrected coefficients approach zero even when observed agreement
is substantial.
AC1 is defined as
\begin{equation}
  \mathrm{AC1} = \frac{p_o - p_e}{1 - p_e},
  \qquad
  p_e = \frac{1}{K-1}\sum_{k=1}^{K} \pi_k(1 - \pi_k),
  \label{eq:ac1}
\end{equation}
where $p_o$ is the observed proportion of agreement,
$K = 3$ is the number of categories, and
$\pi_k = \bigl(\hat{p}_{k,\text{model}} + \hat{p}_{k,\text{gt}}\bigr)/2$
is the mean marginal proportion for category $k$, with $\hat{p}_{k,\text{model}}$
and $\hat{p}_{k,\text{gt}}$ the marginal proportions of class $k$ in the model
predictions and ground truth, respectively.
All metrics---including AC1---are computed globally by pooling all 1{,}589
study--domain cells; domain-level values are not aggregated.
Given the pronounced class imbalance documented in Table~\ref{tab:class-dist}---six
of seven domains have Low Risk as the majority class---AC1 is more appropriate than
Cohen's $\kappa$ as the primary agreement metric for this task, since $\kappa$
approaches zero under heavy prevalence skew even when observed agreement is
substantial (the prevalence paradox; \citealt{Feinstein1990}).
Importantly, the ground-truth annotations used here reflect published Cochrane
consensus judgements rather than a second independent annotator; AC1 therefore
quantifies model-to-published-standard agreement, not inter-rater reliability in
the classical sense.

\textbf{Error direction (LR-anchor framework).}
Because the three RoB classes are not ordinally exchangeable---Low Risk acts as
the ``safe'' default that passes a study through without additional scrutiny,
while High Risk and Unclear Risk both trigger closer examination---we classify
non-diagonal errors relative to Low Risk.
A \emph{lenient} error occurs when the model assigns Low Risk and the ground
truth is High Risk or Unclear Risk (false pass).
A \emph{cautious} error occurs when the model assigns High Risk or Unclear Risk
and the ground truth is Low Risk (false flag).
A \emph{lateral} error occurs when the model confuses High Risk with Unclear Risk
or vice versa.
Because both non-Low-Risk categories already trigger reviewer scrutiny,
lateral errors do not affect the binary pass/no-pass decision; their
consequence is second-order, affecting the precision of the label
(definite bias vs.\ insufficient information) rather than whether the
study receives attention.
This LR-anchor decomposition distinguishes the two first-order error
directions: lenient errors risk including potentially biased studies
without scrutiny (inflating evidence-base credibility), while cautious
errors risk discrediting studies with genuinely low risk of bias
(reducing the available evidence base).
Both carry clinical costs in opposite directions; their relative severity
depends on review context and the scarcity of eligible studies.

\subsection{Consensus Failures and Model Complementarity}
\label{sec:methods-complementarity}

For each study--domain cell, we recorded whether at least one of the eight models
matched the ground truth.
Cells where \emph{no} model was correct are termed \emph{consensus failures} and
represent cases in which all evaluated models failed simultaneously, regardless of
which model was consulted.
The \emph{any-correct rate} (proportion of cells where at least one model
succeeded) characterises the degree of complementarity among models: a high
any-correct rate alongside low individual-model accuracy indicates that different
models fail on different cells, yielding non-redundant coverage across the
ensemble.

\subsection{Stratified Analyses}
\label{sec:methods-strata}

All primary metrics were recomputed within subgroups defined by intervention type
(four categories) and clinical domain (24 Cochrane Review Groups, multi-label).
For clinical domain, studies were included in every group to which they belonged,
so group sample sizes sum to more than 227.

\subsection{Statistical Inference}
\label{sec:methods-inference-stats}
All analyses are descriptive and exploratory.
Confidence intervals are reported to characterise estimation uncertainty and
facilitate comparison across models and conditions; they are not interpreted
as the outcome of pre-specified confirmatory hypothesis tests.
All 95\% CIs were computed by study-clustered bootstrap (5{,}000 iterations):
the 227 studies were resampled with replacement, all seven domain cells of each
sampled study were retained, and the metric of interest was recomputed on the
resulting pseudo-sample; the 2.5th and 97.5th percentiles of the bootstrap
distribution define the interval.
Resampling at the study level preserves within-study correlation across domains
(see Supporting Information~S7 for a comparison with Wilson score intervals).
No familywise error rate correction was applied: the number of comparisons
across models, domains, classes, and inference modes is large and the outcomes
are correlated, making conventional multiplicity corrections either too
conservative (Bonferroni) or difficult to calibrate without a pre-specified
primary endpoint.
Where 95\% bootstrap CIs exclude zero, this is noted descriptively
(``CI excludes zero''); the 95\% level serves as a conventional descriptive
reference threshold and carries no confirmatory implication in this exploratory
context---it marks an unadjusted signal that may warrant further investigation,
not a principled significance claim.

\end{document}
